\makeabstract
    {
    Structure Determination of 2-Chloro-3,3,3-Trifluoropropene and its Complex with Argon via Microwave Rotational Spectroscopy
    }
    {
    Amory Maxey\textsuperscript{1}, Kidane Pincus\textsuperscript{1}, Helen O. Leung\textsuperscript{1}, Mark D. Marshall\textsuperscript{1}
    }
    {
    \textsuperscript{1} Department of Chemistry, Amherst College, Amherst, MA
    }
    {
    Intermolecular forces are a fundamental aspect of chemistry at all levels, and their nuances can be elucidated via microwave rotational spectroscopy. Upon absorption or emission of specific wavelengths of light, molecules transition between quantized rotational states, which are recorded; the spectrum and the resulting structural information are then used to determine the precise structure of the molecule or complex. By applying these techniques to 2-chloro-3,3,3-trifluoropropene and its complex with argon, we further the ongoing halopropene research in the Leung-Marshall lab.
    }
    {
    A 2/3\% sample of 2-chloro-3,3,3-trifluoropropene mixed with argon was analyzed using chirped-pulse Fourier transform microwave spectroscopy from 2 to 18 GHz. Ab initio calculations were performed using the MP2 level of theory using Gaussian16 software to predict rotational constants, which were then used in Western{\textquoteright}s PGopher software to obtain experimental constants. Bond lengths and angles were then calculated using Kisiel{\textquoteright}s STRFIT, yielding cartesian coordinates for the atoms.
    }
    {
    A potential energy scan of the monomer identified a lowest-energy conformation, and the transitions of 5 isotopologues were assigned in the experimental spectrum, allowing determination of 5 sets of rotational constants. Using STRFIT, 6 geometric parameters were fitted to the 15 rotational constants to obtain an experimental structure for the monomer. Scans of the argon complex yielded one minimum, which was identified and fitted on the experimental spectrum, along with its Cl-37 isotopologue. The 6 new constants and the results from the monomer fit were then used to determine the position of the argon atom in the complex, using the same methods.
    }
    {
    The theoretical values proved to be a close prediction for our monomer and argon complex as the differences between predicted and experimental rotational constants were small. We identified the energy minima and obtained experimental structures for the monomer and its complex with argon. Next steps include performing similar analytical techniques on the monomer complex with acetylene.
    }

    \newpage
    